\setcounter{chapter}{0}
\chapter{Topological Spaces}
Non-standard examples of topological spaces to check intuition
\begin{itemize}
    \item (Discrete metric) : $ \Tau \mathcal{T} = \mathcal{P}(X)$
    \item (Indiscrete metric) : $\mathcal{T} = \{\phi, X\} $
    \item $X = \{1, 2, 3\} , \mathcal{T}=\{\phi, \{1\} , \{1, 2\}, \{1, 2, 3\}  \} $
    \item (Co-finite topology) Let $X$ be an infinite set. 
        \[
    \mathcal{T} = \{\phi\} \cup \{U \subseteq X : X \setminus U \text{ is finite }  \} 
        \] 
\end{itemize}
The indiscrete topology is always the coursest topology, and the discrete topology is always the finest, i.e. 
\[
    \{\phi, X\} \subseteq \mathcal{T} \subseteq \mathcal{P}(X)
\] 
\begin{remark}
    We will see many implications of the following form:
    \[
    \text{TOPOLOGICAL PROPERTY} \implies \text{PROPERTY IN TERMS OF SEQUENCES}
    \] 
    In metric spaces, the converse may sometimes be true. In a general topological setting the converse implication is never true.
\end{remark}

\textbf{CLOSURES}
We don't think about closures the way we define them. Instead we think about them in terms of the following property:
\begin{prop}
    $$\overline{A} = \{x \in  X : \text{ for every open } U \subseteq X \text{ with }x \in  U, U \cap A \neq  \O \} $$
\end{prop}
